\chapter{Quiénes aparecen}
\label{cap:personas}

Este apéndice reúne a las personas que el relevamiento documental encontró más de una vez en el departamento, y a unas pocas de aparición única que otras trayectorias necesitan para leerse, con el cargo o la posición que ocupaban en cada aparición y la fecha del
acto que la registra.

No es un índice onomástico ni una lista de responsables. Es una \textbf{herramienta de
trabajo}, y responde a una constatación práctica del capítulo \ref{cap:fincas}: en este territorio, quien
quiera reconstruir el dominio hará bien en empezar por las nóminas de funcionarios, porque
durante la primera mitad del siglo el conjunto de las personas que administraban el
departamento y el de las que eran dueñas de su tierra se superponían casi por completo.

\begin{cautela}
\textbf{Todas las entradas son coincidencias de nombre y apellido en documentos públicos, no
identificaciones registrales.} En un departamento que hasta 1950 tuvo unos pocos miles de
habitantes, los apellidos se repiten, las sucesiones conservan el nombre del causante durante
décadas y padres e hijos comparten nombre de pila.

Ninguna entrada de esta lista prueba que las apariciones correspondan a la misma persona.
Cada una indica dónde buscar para averiguarlo: el Registro Civil, el Registro de Personas
Jurídicas, los padrones electorales y los folios registrales que el apéndice \ref{cap:pedidos} detalla.

Se incluye porque la alternativa --- guardarse las coincidencias --- sería peor: dejaría al
lector sin la posibilidad de verificar el fundamento de la tercera tesis de este libro.
\end{cautela}

\section*{Primera mitad del siglo}

\begin{personas}
\item[{\textbf{Augusto Regis\index{Regis, Augusto}}}] 1908, suplente de los Jurados de Reclamos sobre patentes generales para 1909 (decreto del Ministerio de Hacienda del 12 de noviembre de 1908, B.O.\ Nº 10, 18/11/1908, p.\ 54), donde el nombre se imprime «Augusto Rejis». Febrero de 1909, juez de paz suplente del departamento; renuncia el 1º de marzo. Mayo de 1909, designado miembro de la Comisión Municipal en reemplazo del doctor Rafael Usandivaras ---el decreto no lleva fecha impresa; la edición que lo publica es del 22 de mayo---; ese mismo año figura como lindero al este de la finca Angosto de Arrieta. Titular del jurado de patentes para 1910, por decreto del 4 de noviembre de 1909 \textit{(fecha tomada del OCR del Boletín, no cotejada en la imagen)}. 1910, integrante de la comisión de siete vecinos para las defensas del río. \textbf{18 de abril de 1922, designado miembro de la Comisión Municipal por el Interventor Nacional} ---en el acto que crea ese cuerpo porque el departamento «\textit{en la actualidad carece de ella}»---, confirmado en la rectificación de junio y \textbf{renunciante el 27 de junio, a los dos meses y nueve días} (decreto 132); el tramo leído termina sin acto que cubra su vacante. \textbf{23 de octubre de 1931, como Comisionado Municipal de La Caldera integra la Comisión ad-hoc del Decreto 14091}, junto con el Director General de Obras Públicas y el Comisionado de Campo Santo, encargada de investigar \textbf{todas las concesiones de agua de los ríos Vaqueros y Wierna} y las ordenanzas municipales que las reglaban (cap.\ \ref{cap:redes}). 1933, Presidente de la Comisión Municipal, renombrado al terminar su período anual; ese mismo año Vialidad le encomienda la limpieza del camino de La Caldera a Tres Cruces. 1935, recibe las dos partidas para el espigón de emergencia. \textbf{1937, figura como lindero norte de la finca La Helvecia y, veinte días después, es nombrado nuevamente Presidente de la Comisión Municipal.} 1938, renombrado Presidente. 1939, certifica ante la Provincia de dónde toma el agua la finca «Abra de Lesser». \textbf{17 de diciembre de 1941, el Decreto 3012 lo reelige} Presidente de la Comisión Municipal «por un nuevo período legal». \textbf{5 de mayo de 1942, el Decreto 3829-G lo nombra Comisario de Policía de La Caldera}. \textbf{11 de mayo de 1943, el Decreto 5905-G lo reelige Presidente de la H.\ Comisión Municipal del Distrito de La Caldera} «por un nuevo período legal», por los arts.\ 178 y 182 de la Constitución provincial ---el decreto lo reelige, no lo nombra: ya venía ejerciendo---. \textbf{Septiembre de 1943, el Decreto 526-G acepta su renuncia al cargo de Comisario de Policía de La Caldera}, y el Decreto 1469-G le reconoce los servicios prestados en él. \textbf{En 1943 el presidente del municipio y el comisario de policía del departamento son la misma persona}, según dos decretos de fecha cierta separados por cuatro meses. \textbf{1944, figura en el Revalúo General con una valuación de \$6.400.} \textbf{Treinta y seis años de presencia documentada, en al menos cinco cargos distintos, y propietario en el departamento.} En 1933, además, una \textbf{Julia Regis de Regis} es nombrada Receptora de Rentas y Expendedora de Guías y Marcas de La Caldera (decreto 16.205 del 3 de mayo, B.O.\ Nº 1.486); el parentesco no está acreditado.

\item[{\textbf{Daniel Linares\index{Linares, Daniel}}}] Mayo de 1912, designado para integrar la Comisión Municipal, incompleta por la renuncia de Augusto Rejis\index{Regis, Augusto} \textit{[sic]}; ese mismo año es lindero de San Alejo y Santa Rufina por el este y por el sur, como propietario de la finca Angostura. 1910, integrante y \textbf{presidente} de la comisión de defensas del río. 1911, rinde las cuentas que arrojan el déficit del 48\%. 1914, jurado titular por Caldera, con Bernardo Austerlitz\index{Austerlitz, Bernardo}, para los reclamos de patentes fijas de 1915 (B.O.\ Nº 527); ese mismo año es lindero del naciente y del sur en los tres avisos de remate de San Alejo, dueño de la finca Angostura, y \textbf{demandado ante el Juez Federal en un juicio de acción negatoria promovido por José Cobach}, en el que se rematan sin base diecinueve vacunos depositados en poder de Antonio Villegas, en La Caldera (B.O.\ Nº 501). 1922, designado miembro de la Comisión Municipal. \textbf{1932, figura en el deslinde de las cinco fincas de Silvano Y.\ Murúa como propietario de la finca y la estancia Los Porongos}. 1934, su apellido identifica una de las casas donde se cometen los hurtos que resuelve la Sala en lo Penal. 1938, lindero sur de la finca Campo Alegre con su finca El Angosto. 1943, propietario de los terrenos que lindan al sur y al este con la Comisaría de Policía del pueblo, sobre cuyo predio el Estado dedujo posesión treintañal

\item[{\textbf{Ángel Fernández\index{Fernández, Ángel}}}] 1908, Presidente Municipal; firma el cuadro de rentas del ejercicio. Enero de 1909, designado miembro de la Comisión Municipal; ese año es comisario auxiliar del partido de San Alejo y Santa Rufina y \textbf{reconoce en juicio haber sido administrador de las fincas San Alejo y Santa Rufina}, en división de condominio de su propia familia. 1910, integrante de la comisión de defensas del río

\item[{\textbf{Silvano Ignacio Murúa\index{Murúa, Silvano Ignacio}}}] Nacido en Salta el 31 de julio de 1882, hijo de Silvano Murúa y de Avelina Cornejo; casado el 24 de diciembre de 1908 con Fanny Ovejero Lacroix; muerto el 31 de diciembre de 1947 \textit{(fechas de la compilación genealógica citada en el apéndice \ref{cap:fuentes})}. \textbf{Febrero de 1909, a los veintiséis años, dos meses después de casarse, renuncia a la Comisión Municipal}; el decreto lo nombra «Silvano J.\ Murúa» y lo reemplaza Abelardo Figueroa\index{Figueroa, Abelardo}. 1909 y 1910, suplente en las mesas receptoras de la sección Caldera. 1910, a los veintisiete, integra la comisión de defensas del río. 1924, a los cuarenta y uno, hace deslindar con su esposa la finca Wierna; ese año un «Silvano Murúa» ---él o su padre, homónimo y todavía vivo--- figura como lindero de una finca de pastoreo en la quebrada de Mojotoro. 1927, candidato a diputado por Caldera de la Unión Provincial, y, como «Silvano I.\ Murúa», hace deslindar el Puesto de las Garrapatas. \textbf{1932, a los cuarenta y nueve, como «Silvano Y.\ Murúa», él y su esposa mandan deslindar cinco fincas, entre ellas «El Durazno o Peñones»}. 1934, el puesto Cañada Seca o Peñones figura «adjudicado al heredero Silvano Murúa». 1936, lindero norte de La Despensa, rematada en la ejecución contra tres herederos Murúa. 1944, a los sesenta y dos, figura en el Revalúo General con una valuación de \$7.600. \textit{Las cuatro grafías ---J., I., Y.\ e Ignacio--- corresponden a la misma persona}

\item[{\textbf{Silvano Murúa\index{Murúa, Silvano}} \textit{(padre)}}] Nacido en Salta el 3 de mayo de 1855; casado en Campo Santo en 1881 con Avelina Cornejo; padre de Silvano Ignacio, de Eulalia ---la «Eulalia Murúa de Cornejo» que en 1934 recibe el puesto Los Matos--- y de Carlos, Eustaquio, Juan Avelino, María Margarita y Víctor \textit{(misma compilación)}. En noviembre de 1914 un «Silvano Murúa», sin segundo nombre, es jurado suplente por Caldera para los reclamos de patentes (B.O.\ Nº 527); puede ser él o su hijo, y el decreto no permite decidirlo. \textbf{En agosto de 1927 ya había muerto}: el edicto del Puesto de las Garrapatas nombra como linderos a «la sucesión de don Silvano Murúa», y los de 1932 y 1933 vuelven a nombrarla. Es la sucesión de la que su hijo es heredero

\item[{\textbf{Abelardo Figueroa\index{Figueroa, Abelardo}}}] 1908, titular de los Jurados de Reclamos sobre patentes generales para 1909. Febrero de 1909, designado miembro de la Comisión Municipal en reemplazo de Silvano Ignacio Murúa\index{Murúa, Silvano Ignacio}; ese año es lindero al naciente de San Alejo y Santa Rufina y vuelve a ser titular del jurado de patentes, para 1910. 1910, integrante de la comisión de defensas del río

\item[{\textbf{Eustaquio Murúa\index{Murúa, Eustaquio}}}] 1885, la Municipalidad le concede treinta litros por segundo del río Wierna para la finca «Sausal o Curuzú» (apéndice \ref{cap:cronologia}). \textbf{De auxiliar de un partido a los tres cargos del departamento en cuatro años.} Abril de 1912, comisario auxiliar del partido de La Despensa, cargo que conserva en 1913; en septiembre, sorteado titular de la segunda mesa receptora de votos del departamento. \textbf{Enero de 1915, nombrado miembro de la comisión municipal} del departamento por decreto del Poder Ejecutivo, con Adolfo Vidal (B.O.\ Nº 540, h.\ 2); en septiembre del mismo año, \textbf{candidato a elector de gobernador por el partido Unión Provincial}, con Nolasco Echenique (Nº 569, h.\ 2). \textbf{1916, tres cargos en cinco meses}: comisario de policía del departamento el 4 de marzo de 1916, en la vacante de Manuel R.\ Cornejo (Nº 596, h.\ 1), cargo que conserva en 1917 (decreto 1101, Nº 639, h.\ 5); expendedor de guías el 16 de marzo, con fianza de Silvano I.\ Murúa\index{Murúa, Silvano Ignacio} por \$500 (Nº 597, h.\ 2); y receptor del impuesto al consumo por La Caldera el 5 de agosto (Nº 616, h.\ 2). Coincidencia de nombre a veintisiete años de distancia, no identidad probada

\item[{\textbf{Salvador Rosa\index{Rosa, Salvador}}}] 1913, miembro de la Comisión Municipal del departamento, con Daniel Linares\index{Linares, Daniel} y Camilo Cáceres. 1936, ejecutante en el juicio en que se remata La Despensa contra Abel, Adolfo y Alberto Murúa (apéndice \ref{cap:dominio}). Es también el propietario de El Nogalito en el revalúo de 1944

\item[{\textbf{Enrique Leconte\index{Leconte, Enrique}}}] 1912, con Pablo Enrique Wernner, pide el deslinde de la finca Los Porongos o San Luis. Abril de 1913, comisario auxiliar del partido de Los Porongos y, siete días después, del de Los Peñones

\item[{\textbf{Benjamín López\index{López, Benjamín}}}] 1909, subcomisario del partido de la estación Mojotoro. Octubre de 1911, encargado \textit{ad honorem} de la oficina del Registro Civil de la estación. 1912, declara en juicio como jefe de la estación del ferrocarril y en diciembre es nombrado comisario del partido de Mojotoro

\item[{\textbf{Rafael Fernández\index{Fernández, Rafael}}}] 1911, juez de paz propietario del departamento, en reemplazo de Carlos Matienzo\index{Matienzo, Carlos}, que renunció. 1912, nombrado otra vez para el ejercicio y firma en Caldera, como juez de paz, el edicto sucesorio de Rafael Peralta

\item[{\textbf{Juan Chuchuy\index{Chuchuy, Juan}}}] 1908, actor en un juicio por cobro de pesos; la sentencia lo da «domiciliado en el Departamento de la Caldera». 1909, propietario de mesa receptora de la sección Caldera. \textbf{1911, comisario auxiliar de policía del partido de Yacones}, y otra vez en 1912

\item[{\textbf{Carlos Matienzo\index{Matienzo, Carlos}}}] 1908, firma con el Presidente Municipal el cuadro de rentas del departamento. Febrero de 1909, juez de paz propietario; en abril preside el acta que nombra los veedores del Código Rural. Julio de 1910, otra vez juez de paz propietario del departamento, con Jacinto Guerrero\index{Guerrero, Jacinto} de suplente; febrero de 1911, lo nombran por tercer año consecutivo y en mayo renuncia. Lo reemplaza Rafael Fernández

\item[{\textbf{Escolástico M.\ Aparicio\index{Aparicio, Escolástico M.}}}] 12 de noviembre de 1908, clasificador de patentes generales de Caldera para 1909 (B.O.\ Nº 10, 18/11/1908, p.\ 54). 23 de diciembre de 1908, comisario de policía del departamento para 1909, ya sin la inicial (B.O.\ Nº 22, 30/12/1908, p.\ 105). Mismo nombre, mismo departamento y seis semanas de diferencia. Otra vez clasificador de patentes de Caldera, para 1910, por decreto del mismo día que ese jurado. Febrero de 1912, otra vez comisario de policía del departamento. Febrero de 1913, renuncia a la comisaría; abril de 1913, \textbf{juez de paz propietario} del departamento, por terna de la comisión municipal; \textbf{9 de mayo de 1914, miembro de la comisión municipal} en la vacante que deja Salvador Rosa\index{Rosa, Salvador} al mudarse a la capital (B.O.\ Nº 485, 1914) ---séptimo año seguido con cargo en el departamento---; en agosto y octubre firma como tal los edictos sucesorios del pueblo, y en noviembre es comisionado para revisar las patentes de las casas de comercio. \textbf{1915: juez de paz y clasificador de patentes en el mismo año.} El 2 de septiembre firma como juez de paz del departamento el edicto de la sucesión de Asunción Espeche González de Wayar (B.O.\ Nº 568, h.\ 4, y Nº 569, h.\ 3), y el 6 de diciembre un decreto lo comisiona para la clasificación general de patentes y el empadronamiento de los capitales en giro de 1916 (Nº 584, h.\ 2). \textbf{1916: renuncia y vuelve.} En febrero un decreto acepta su renuncia como miembro de la comisión municipal y nombra en su lugar a Andrés Marín (B.O.\ Nº 593, h.\ 1); en abril, «de acuerdo con la terna elevada por la Comisión municipal de la Caldera» ---el mismo cuerpo que acababa de dejar---, es nombrado \textbf{Juez de Paz propietario} del departamento para el ejercicio del año (Nº 601, h.\ 3). \textbf{1911}, comisionado del empadronamiento de los capitales en giro y de la clasificación de patentes para 1912, por decreto del Ministerio de Hacienda del 9 de noviembre (B.O.\ Nº 295, h.\ 2): es el año que cerraba la serie y que el barrido de 1911 registra. \textbf{1917}, juez de paz propietario del departamento para el ejercicio del año, por terna de la comisión municipal renovada diecinueve días antes (decreto 1085; B.O.\ Nº 638, h.\ 7), con Casimiro Cáceres de suplente. \textbf{Y la serie termina en 1918, y termina por un traslado.} Un decreto del Interventor Nacional, publicado en el B.O.\ Nº 704 (h.\ 11), traslada personal de la Receptoría de Rentas: \textbf{Escolástico Aparicio, de Caldera a Campo Santo}, y Pedro Pérez del Busto en sentido inverso. \textbf{Once años seguidos en uno u otro cargo, de 1908 a 1918, y en seis de ellos ---1908, 1909, 1910, 1911, 1913 y 1915--- el cargo es el que clasifica o empadrona a los contribuyentes del departamento} (cap.\ \ref{cap:siglo}). \textbf{No lo saca una elección ni una renuncia: lo saca la Intervención, y lo saca del departamento.} \textbf{Y vuelve en 1921}: el decreto 1669 del 7 de julio lo nombra \textbf{Juez de Paz Suplente} del departamento, otra vez por terna de la Comisión Municipal, y el 1816 del 12 de septiembre acepta su renuncia «\textit{en atención a los motivos en que la funda}», que no transcribe. \textbf{Dos meses y cinco días, y sin reemplazante publicado.} Con eso, doce apariciones en catorce años El Boletín lo imprime «Escolástico M.», «Escolástico C.» y «Escolástico» a secas. Como comisario del departamento propone los comisarios de partido de 1909, el de la estación Mojotoro y el de La Despensa. \textbf{1910: cuatro cargos.} Marzo, miembro de la Comisión Municipal en reemplazo de Abelardo Figueroa\index{Figueroa, Abelardo}; abril, comisario de policía del departamento; diciembre, comisionado de empadronamiento y clasificación de patentes y, en el mismo mes, otra vez comisario y otra vez miembro de la Comisión Municipal. El Boletín lo imprime «Escolástico M.\ Aparicio», «Escolástico Aparicio» y «Escolastino Aparicio»

\item[{\textbf{Jacinto Guerrero\index{Guerrero, Jacinto}}}] Enero de 1909, miembro de la Comisión Municipal; febrero, comisario suplente de policía del departamento; marzo, juez de paz suplente en reemplazo de Augusto Regis\index{Regis, Augusto}; abril, renuncia a la comisión municipal. Firma además, como vecino propietario, el acta de los veedores del Código Rural. 1910, y otra vez en 1912, comisario suplente del departamento y juez de paz suplente, los dos cargos a la vez, y integrante de la comisión de defensas del río

\item[{\textbf{Dermidio Arancibia}}] \textbf{1918}, con Miguel A.\ Sosa, peticionante de la mina de plomo «20 de Febrero» en la finca Cerro Nevado y denunciante del despueble de la mina «Sara», las dos del departamento (B.O.\ Nº 711 y Nº 716). \textbf{1920}, candidato a diputado por La Caldera proclamado por la Unión Cívica Radical (B.O.\ Nº 796). \textbf{Coincidencia de nombre, no identificación registral}; si es la misma persona, es el primer caso de esta serie en que quien pide una concesión sobre el suelo del departamento se presenta a representarlo

\item[{\textbf{Liborio Guerrero\index{Guerrero, Liborio}}}] 1909, comisario auxiliar del partido de Peñones y Chalchanio; \textbf{1918, el mismo cargo y el mismo par de partidos}, por el decreto 1.564, nueve años después. \textbf{1919, comisario del partido del Monte} por el decreto 213 de marzo, y en mayo renuncia a «\textit{los partidos de Monte y Chalchalnio}» (decreto 337): lo reemplaza Héctor Echanique. \textbf{Tres apariciones en once años, siempre en el mismo par de partidos del este del departamento.} 1932, sus herederos figuran como linderos en el deslinde de las cinco fincas de los Murúa

\item[{\textbf{Genaro Molina\index{Molina, Genaro}}}] 1909, comisario auxiliar del partido de La Calderilla. 1937, figura como lindero sur de la finca La Helvecia, antes Calderilla

\item[{\textbf{Augusto Regis\index{Regis, Augusto}}}] \textbf{18 de abril de 1922}, designado a la Comisión Municipal por el Interventor Nacional, en el acto que crea ese cuerpo porque el departamento «\textit{en la actualidad carece de ella}»; confirmado en la rectificación de junio y \textbf{renunciante el 27 de junio}, a los dos meses y nueve días. \textbf{18 de enero de 1923}, el decreto 620 acepta su \textbf{segunda renuncia} al mismo cuerpo. \textbf{1937 a 1940}, \textbf{presidente de la Comisión Municipal del distrito}, reelecto por decreto cuatro años seguidos ---«un nuevo período legal de funciones»---. \textbf{1937}, lindero norte de «La Helvecia» antes «Calderilla» en su edicto de deslinde. \textbf{1939}, firma el certificado sobre la provisión de agua que un propietario presenta ante la Provincia (B.O.\ Nº 1778, h.\ 15). \textbf{1940}, el Boletín lo identifica como «\textit{presidente de la comisión municipal del distrito de la Caldera, \textbf{diputado provincial}}» (Nº 1875, h.\ 2). \textbf{Dieciocho años en el mismo cuerpo, con dos salidas en el medio}

\item[{\textbf{Carlos F.\ Matienzo\index{Matienzo, Carlos}}}] \textit{(Coincidencia de nombre entre 1909 y 1922, no identificación registral.)} \textbf{1909}, juez de paz propietario del departamento, y firma además como secretario del cuadro de rentas municipales de 1908 que el Boletín publica en febrero. \textbf{12 de junio de 1922}, \textbf{juez de paz propietario} del departamento por terna de la Comisión Municipal recién creada, con Vicente Vivas de suplente (decreto 105; B.O.\ Nº 912, h.\ 3). \textbf{27 de junio de 1922}, \textbf{encargado del Registro Civil} del mismo departamento (decreto 135; Nº 914, h.\ 4): \textbf{quince días entre un nombramiento y el otro}, y el Boletín no dice si los cargos son compatibles. Es el cuarto encargado del Registro Civil del departamento en cuatro meses (cap.\ \ref{cap:siglo})

\item[{\textbf{Francisco S.\ Urquiza\index{Urquiza, Francisco S.}}}] \textit{(«Francisco Solano» en la nómina de presidentes de comicio insaculados para 1928, circuito 6; la identidad es probable y no está acreditada.)} 1921, hace deslindar judicialmente las fincas Cerro de Buena Vista y Sausal o Curuzú. 1924, obtiene concesión de cien litros por segundo del río Wierna para regar Cerro de Buena Vista. \textbf{Las dos fincas que hizo deslindar en 1921 sacan agua del mismo río}: el Decreto 5933 de 1942 (cap.\ \ref{cap:redes}) reconoce que «Sausal o Curuzú», distrito de Vaqueros, tiene \textbf{treinta litros por segundo del Wierna} por concesión que la Municipalidad de La Caldera otorgó el 4 de junio de \textbf{1885} a Eustaquio Murúa. Sumadas, las dos fincas llevan ciento treinta litros por segundo del Wierna. \textbf{Y el edicto de 1921 dice con qué título lo hizo}: el auto del juez Daniel Etcheverry manda deslindar ambas fincas, en el partido de Vaqueros, «\textit{de propiedad de don Francisco S.\ Urquiza}», con agrimensor Héctor Chiostri y sobre los títulos de fs.\ 33 y 40 (B.O.\ Nº 843 del 28 de enero de 1921). \textbf{Era dueño de las dos.} El decreto de 1942 no lo nombra ---da la cadena como Murúa, «antecesor en la posesión», y Guillermo Schaefer como propietario, sin titular intermedio---, de modo que entre 1921 y 1942 «Sausal o Curuzú» cambió de manos y el reconocimiento provincial omitió el eslabón. \textit{Cuándo y cómo se transfirió se contesta con el folio.} \textbf{Hasta abril de 1928, miembro de la Comisión Municipal de La Caldera}, cuyo informe sustenta ese mes el rechazo del reclamo de Campo Santo contra su propia concesión; renuncia nueve días después. \textbf{Agosto de 1928, sometido a proceso penal «por hurto y uso indebido de agua», con prisión preventiva}, por denuncia de Juan A.\ Bianchi\index{Bianchi, Juan A.}, propietario de la finca vecina. 1930, sobreseído por prescripción. 1931, la Corte de Justicia anula la revocación. 1935, lindero este de la finca «Vaqueros o Entre Ríos». 1936, la Corte rechaza su demanda de daños por \$20.000. \textbf{1944, figura en el Revalúo General como el mayor de los cuarenta y siete titulares publicados}, con una valuación de \$52.400 --- casi el quince por ciento de todo el valor publicado

\item[{\textbf{Bernardo Austerlitz\index{Austerlitz, Bernardo}}}] 1914, jurado titular por Caldera, con Daniel Linares\index{Linares, Daniel}, para los reclamos de patentes fijas de 1915 (B.O.\ Nº 527): es su primera aparición en el archivo. 1925, propietario del suelo de la finca San Alejo y Santa Rufina, según dos pedimentos de cateo de dos mil hectáreas cada uno; domiciliado en Alberdi 374 de la capital. \textbf{1927, senador electo por el departamento de Caldera}. 1938, su heredera Carolina Austerlitz de Aráoz figura como lindero oeste de Campo Alegre con el inmueble Santa Rufina

\item[{\textbf{Hermann Pfister\index{Pfister, Hermann}}}] 1928, designado miembro de la Comisión Municipal que reemplaza a la declarada en acefalía (Decreto 9306, junio). 1932, perito designado en el deslinde de las cinco fincas de Silvano Ignacio Murúa y su esposa. 1934, perito en el deslinde del Puesto Despensa. \textbf{1937, solicitante del deslinde de La Helvecia}. En el mismo período figura como inspector auxiliar de minas. \textbf{1945, promueve juicio de posesión treintañal} sobre un terreno «con todo lo plantado» en La Calderilla, de 108,28 m de frente por 2.395 m de fondo, cuyo lindero \textbf{Oeste es el río de La Caldera} y cuyo lindero Norte es, según el propio edicto, «propiedad del Sr.\ Pfister» ---sin nombre de pila: la coincidencia se anota, el parentesco ni la identidad se afirman---. Auto del 3 de julio de 1945, Juzgado Civil de Tercera Nominación, Dr.\ \textbf{Alberto E.\ Austerlitz\index{Austerlitz, Alberto E.}}; edicto Nº 1214, treinta ediciones consecutivas (B.O.\ Nº 2.403 a 2.432). \textbf{El perito que midió las fincas ajenas aparece trece años después adquiriendo por prescripción un inmueble ribereño en el mismo departamento}

\item[{\textbf{Juan Larrán\index{Larrán, Juan}}}] 1925, cateo sobre la finca Las Nieves (expte.\ 1168 C). 1930, el 30 de abril se registra su manifestación de las minas Caldera Nº 1, 2 y 3 en la finca San José (expte.\ 44-L); el 15 de octubre se le concede, con Alberto Jaime Harrison, un cateo sobre Las Nieves pedido en 1929 (expte.\ 38-H). 1933, titular de esas minas en el padrón minero. 1941, las minas se rematan por falta de pago del canon, junto con «La Colorada» (cobre, La Poma), que tenía con Alejandro Bonari. En 1941 un \textbf{Oscar Larrán Sierra} es nombrado Inspector General del Departamento Provincial del Trabajo (decreto 1528 G, B.O.\ Nº 1.898); el parentesco no está acreditado

\item[{\textbf{Rafael Delgado\index{Delgado, Rafael}}}] 1936, juez de paz del distrito municipal de La Caldera. \textbf{1936, ejecutado en el remate sin base del aserradero del pueblo} --- sierra sin fin, dos circulares, motor a vapor, garlopa y dieciséis bueyes --- a instancia de una casa comercial de la capital

\item[{\textbf{Ramón Juárez\index{Juárez, Ramón}}}] 1930, Encargado del Registro Civil de La Caldera (decreto 12.406 del 13 de octubre, B.O.\ Nº 1.346). 1932--1937, comisario de policía de La Caldera. 1937, \textbf{Ramón I.\ Juárez} es el ejecutado en el remate de la finca Chalchanio. La inicial del segundo nombre aparece sólo en el edicto

\item[{\textbf{Juan A.\ Bianchi\index{Bianchi, Juan A.}}}] 1928, denunciante en la causa penal por hurto y uso indebido de agua contra Francisco S.\ Urquiza\index{Urquiza, Francisco S.}; \textbf{propietario de la finca «Vaqueros» o «Entre Ríos»}, que en 1935 el Banco Provincial remató en la ejecución contra Ramón Bracheri, ya otro titular. En 1929, una «señora de Bianchi» figura como lindero norte de la casa quinta de cinco hectáreas rematada en Vaqueros

\item[{\textbf{Carlos Serrey\index{Serrey, Carlos}}}] 1939, propietario de los terrenos que lindan por el norte y el oeste con el inmueble sobre el que la Curia Eclesiástica promueve posesión treintañal en el pueblo de La Caldera. Su sucesión, resuelta a comienzos de los años noventa, es identificada por la literatura académica como origen del mercado de suelo de Vaqueros; su sucesión donó el terreno del edificio municipal de ese municipio; en 1986 se expropió a Elena Serrey\index{Serrey, Elena} de González Bonorino un inmueble de La Caldera para viviendas familiares; en 1985 la Ley 6334 había expropiado 2 ha 1.736,80 m² de la Finca Vaqueros, a nombre de Manuel Serrey y otros, para viviendas; y en octubre de 1986 un decreto de necesidad y urgencia, convertido luego en la Ley 6464, expropió 3.710,08 m² de la misma finca, de la sucesión testamentaria de Carlos Serrey, para la planta potabilizadora de Vaqueros

\item[{\textbf{Teófilo Reyes\index{Reyes, Teófilo}}}] 9 de mayo de 1928, Encargado del Registro Civil de La Caldera en reemplazo de Ramón Balboa (Decreto 8002, B.O.\ Nº 1.222). 1943, Juez de Paz Propietario de La Caldera; diligencia el oficio del juicio de posesión treintañal que el Estado promovió sobre el predio de la Comisaría. Comparte apellido con el \textbf{Teodoro Reyes} que en 1909 figura como comisario auxiliar del partido de Los Porongos

\item[{\textbf{Adeodato Aybar\index{Aybar, Adeodato}}}] 1929, catastrador; termina el levantamiento de los departamentos de Campo Santo, Caldera y La Viña

\item[{\textbf{Adolfo Vidal} y \textbf{Carlos Murúa\index{Murúa, Carlos}}}] 1908, titulares de los Jurados de Reclamos sobre patentes generales para 1909 (decreto del Ministerio de Hacienda del 12 de noviembre de 1908, B.O.\ Nº 10, 18/11/1908, p.\ 54), junto con Abelardo Figueroa\index{Figueroa, Abelardo}; el decreto imprime «Cárlos». Abril de 1909, Adolfo Vidal es designado miembro de la Comisión Municipal en reemplazo de Jacinto Guerrero\index{Guerrero, Jacinto}, y vuelve a ser titular del jurado de patentes, para 1910. 1909, Carlos Murúa es sorteado para la mesa de la sección Caldera junto con Silvano Murúa. \textit{No es Carlos Murúa Cornejo, hermano de Silvano Ignacio, que había nacido en 1895 y tenía trece años}

\item[{\textbf{José Miguel Rauch\index{Rauch, José Miguel}}}] 2 de junio de 1910, comisario del partido de La Calderilla, nombrado de acuerdo con la propuesta del comisario de policía del departamento (B.O.\ Nº 161, 1910). Octubre de 1910, damnificado en una causa por hurto de ganado sustraído de su finca en La Calderilla (B.O.\ Nº 200, 1910): el comisario del partido era propietario en el partido que vigilaba (cap.\ \ref{cap:fincas}). \textit{El apellido vuelve dos veces más y este libro no establece vínculo alguno}: un \textbf{Jorge Rauch} integra en 1922 la Comisión Municipal que designa el Interventor Nacional (cap.\ \ref{cap:siglo}), y una \textbf{Celina Z.\ de Rauchi} figura en el Revalúo General de diciembre de 1944 bajo la partida 104, con \$36.500: la segunda valuación más alta del departamento (caps.\ \ref{cap:fincas} y \ref{cap:expropiacion})

\end{personas}

\section*{Segunda mitad del siglo y presente}

\begin{personas}
\item[{\textbf{Arturo R.\ Fernández\index{Fernández, Arturo R.}}}] 1983, Presidente de la Comisión Municipal de La Caldera; firma la Ordenanza 047 que rebaja tasas a jubilados, sancionada y promulgada por él mismo en virtud del Decreto 877/82

\item[{\textbf{Esteban Mogro\index{Mogro, Esteban}}}] 1983, Secretario Administrativo de la Comisión Municipal; firma la Ordenanza 047 en junio. 1983, renuncia en noviembre al cargo de \textbf{Juez de Paz Titular} de la localidad. Pueden ser cargos sucesivos o simultáneos, o dos homónimos

\item[{\textbf{Héctor Miguel Calabró\index{Calabró, Héctor Miguel}}}] Intendente de La Caldera entre 1999 y 2011, en tres mandatos consecutivos (fuentes periodísticas coincidentes; los Decretos 3318/08 y 2508/09 lo nombran como intendente en funciones). 2007, firma como intendente el acta del Servicio de Información Catastral (Decreto 1639/07). 2008, presidente de la Cooperadora Asistencial de La Caldera por designación provincial. 2011, diputado provincial (acta de proclamación). Diciembre de 2015, secretario de Relaciones Institucionales del Ministerio de Gobierno (Decreto 89/15); 23 de noviembre de 2017, renuncia a ese cargo (Decreto 1608/17) y es designado Secretario de Asuntos Municipales de la Provincia (Decreto 1634/17), cargo que ejerce en 2018 (Decisión Administrativa 592/18) y, según la prensa, hasta diciembre de 2019. 8 de noviembre de 2018, mientras es Secretario de Asuntos Municipales, designado gerente suplente de Jardín Celestial S.R.L. (acta publicada en enero de 2020; sin constancia posterior de cese ni de continuidad). 2021, senador provincial por el departamento, hasta diciembre de 2025. 2023, presenta el proyecto de Reserva Natural Abra de Santa Laura, que caduca sin tratamiento en 2025. 2026, la Provincia acepta dos donaciones con cargo que él le hace: una motocicleta para la Comisaría 15 de Vaqueros (Decreto 59/26) y otra para la Subcomisaría de La Caldera (Decreto 225/26), ambas para patrullaje

\item[{\textbf{Carlos Miguel Calabró\index{Calabró, Carlos Miguel}}}] 2010, socio fundador de Jardín Celestial S.R.L.\ con el 30\% del capital; diecinueve años, estudiante, domiciliado en La Caldera. 2014, tras el aumento de capital, pasa de 540 a 1.635 de las 5.450 cuotas: conserva exactamente el 30\%. \textit{El vínculo familiar con el anterior no está acreditado con partida}

\item[{\textbf{Héctor D'Andrea\index{D'Andrea, Héctor}}}] Socio y gerente titular de Jardín Celestial S.R.L.; 45,01\% del capital desde 2014; domiciliado en Finca Getsemaní, La Caldera. Copropietario de la matrícula 6.558, sobre la que se aprobó el plano del loteo El Durazno

\item[{\textbf{Magdalena Vidal\index{Vidal, Magdalena}}}] 2015, gestiona como \textbf{co-titular registral del catastro 4366} la concesión de agua para cuatrocientos lotes del loteo El Durazno

\item[{\textbf{Silvia F.\ Santamaría\index{Santamaría, Silvia F.}}}] 2009 y 2011, jefa de subprograma; firma avisos del organismo. 2014--2015, Abogada Jefe del Programa Fiscalización y Control; firma edictos de concesión de agua, entre ellos los del Catastro 1700, El Durazno y La Misión III. De 2019 hasta al menos julio de 2025, Jefa de Programa Jurídico. Desde diciembre de 2025, Directora Legal y Técnica de la Subsecretaría de Recursos Hídricos; sigue firmando edictos del organismo

\item[{\textbf{María Silvina Abilés\index{Abilés, María Silvina}}}] Senadora provincial por La Caldera, en 2021 disputó sin éxito la reelección frente a Calabró (cap.\ \ref{cap:politica}). Desde el 10 de diciembre de 2023, Secretaria de Tierras y Bienes del Estado (Decreto 64/23, cap.\ \ref{cap:tierrafiscal}); sin constancia de reemplazo en 2026. Un medio la señala como copropietaria de un loteo del departamento; este libro no lo afirma (cap.\ \ref{cap:loteo})

\item[{\textbf{Sergio Camacho\index{Camacho, Sergio}}}] Ministro de Infraestructura al menos entre 2021 y 2025 (cap.\ \ref{cap:tierrafiscal}); en 2021 firma la Resolución 646D/21, que aprueba el plano del loteo El Durazno (cap.\ \ref{cap:loteo}). Desde el 15 de diciembre de 2025, Jefe de Gabinete de Ministros, que absorbe las áreas del ex Ministerio de Infraestructura; en 2026 convoca la reunión de enero en la que se acepta el plan de represas de retardo (cap.\ \ref{cap:amparo})

\item[{\textbf{Matías J.\ Brogin\index{Brogin, Matías J.}}}] 2012--2016, Asesor Legal de la Secretaría de Recursos Hídricos; firma los edictos de la serie de línea de ribera del departamento ---salvo el de la Resolución 023/14, que firma el secretario Fuertes---, incluidas las tres determinaciones que no publicaron coordenadas

\item[{\textbf{Jorge G.\ Torres\index{Torres, Jorge G.}}}] 2014, designado para determinar la línea de ribera del río La Caldera a lo largo del pueblo; la misma resolución la determina, con coordenadas publicadas en su anexo. 2015, integra dos comisiones técnicas del departamento. 2016, la Resolución 360/16 aprueba la determinación del río Wierna hecha por la comisión que había integrado en 2015 (Res.\ 319/15). 2016, firma además como director general de Planificación Hídrica la determinación «por Administración» del arroyo Quintín (Res.\ 389/16), sobre el mismo tramo de una comisión que integró ese año. 2025, Director General de Cuencas y Política Hídrica. Autor de la presentación institucional sobre línea de ribera que este libro usa como fuente del procedimiento

\item[{\textbf{Pablo Vicente Balverdi}}] 2012, ingeniero en construcciones de la comisión que determina la línea del río Castellanos en la matrícula 4027 (Res.\ 94/12)

\item[{\textbf{Enrique Chalabe}}] 2016, licenciado en ciencias geológicas en la comisión del arroyo Chaile (Res.\ 290/16); 2025, autorizado a diligenciar la publicación de la rectificación de la línea del río Vaqueros sobre la matrícula 216 (Res.\ 130/2025)

\item[{\textbf{Rolando Chávez}, \textbf{Víctor Zacarías Pérez} y \textbf{Héctor Saravia Navamuel}}] 2016, comisión técnica del río Vaqueros, catastro 40 (Res.\ 147/16): ingeniero en construcciones, ingeniero hidráulico y geólogo

\item[{\textbf{Marcelo Ricardo Toigo} y \textbf{José A.\ Medina}}] 2011, ingeniero civil y geólogo de la comisión técnica del arroyo Gallinato en Chacras del Gallinato. Toigo integra en 2016 la del arroyo Chaile (Res.\ 290/16)

\item[{\textbf{Antonio Alejandro Forns\index{Forns, Antonio Alejandro}}}] 2013, integra la comisión técnica del arroyo Quintín y la del río La Caldera y arroyo El Manzano; 2015, la del río Trancas o Wierna en Finca Vaqueros (Res.\ 273/15); 2025, autorizado a diligenciar la publicación de la determinación del arroyo Quintín que cierra la comisión de 2013 (Res.\ 142/2025)

\item[{\textbf{Mauricio Romero Leal}}] 2022, firma el informe de mapas de la cuenca. 2023, secretario de Recursos Hídricos, destinatario de la nota del 20 de enero. Febrero de 2026, designado Subsecretario de Recursos Hídricos (Decreto 52/26)

\item[{\textbf{David Le Favi}}] 2023, director de Fiscalización y Control del organismo hídrico; firma la nota del 20 de enero y el Dictamen 129/2023. Diciembre de 2025, coordinador de la Subsecretaría. Febrero de 2026, representa al organismo en la audiencia de aprobación del Plan

\item[{\textbf{Rodrigo Matías Páez}}] 2022, como subsecretario de la Municipalidad de La Caldera, firma la nota 83/22 que eleva a Fiscalía de Estado el registro de loteos de la microcuenca (cap.\ \ref{cap:amparo}). Responsable municipal de Medio Ambiente, Tierra y Hábitat en las gestiones ante la Secretaría de Tierras y Bienes del Estado (cap.\ \ref{cap:tierrafiscal})

\item[{\textbf{Dal Borgo} (César Domingo; Dal Borgo Construcciones S.R.L.)}] 2014, la sociedad es adjudicataria de la protección de las cámaras del Acueducto Norte. 2018--2020, César Domingo Dal Borgo figura como concesionario de la cantera «La Vaquera» y de una servidumbre de acopio, con informes no aprobados. 2022, la sociedad es oferente en la licitación del puente sobre el río Vaqueros. \textit{La relación entre la persona y la sociedad no está acreditada}

\item[{\textbf{Daniel Roberto Moreno Ovalle\index{Moreno Ovalle, Daniel Roberto}}}] Intendente de Vaqueros electo en 2011 (Salta Somos Todos), 2015 (Partido de la Victoria), 2019 y 2023. En mayo de 2025 es electo senador provincial por el departamento; renuncia a la intendencia en diciembre de 2025 para asumir la banca. En la elección de 2015, \textbf{Mario Enrique Moreno Ovalle} --- con \textit{número de documento consecutivo}, según el acta --- fue electo concejal de la ciudad de Salta por la misma fuerza

\item[{\textbf{Luis Gerardo Mendaña\index{Mendaña, Luis Gerardo}}}] Intendente de La Caldera electo en 2011. Pide a Recursos Hídricos la línea de ribera del río a lo largo del pueblo (Res.\ 023/14). En marzo de 2014 firma, como intendente, el convenio de licencias de conducir que el Decreto 1433/15 aprueba en 2015; en 2025 respalda como diputado la candidatura de Calabró

\item[{\textbf{Daniel Alejandro Escalera\index{Escalera, Daniel Alejandro}}}] Intendente de La Caldera electo en 2015; en ejercicio durante el aluvión de diciembre de 2018

\item[{\textbf{Diego Ramiro Sumbay\index{Sumbay, Diego Ramiro}}}] Intendente de La Caldera desde 2019, reelecto en 2023. En 2015, \textbf{Alfio Edgardo Sumbay} había sido electo concejal del municipio por otra fuerza política

\end{personas}

\section*{Cómo se usa este apéndice}

Cada fila indica un lugar donde buscar, no una conclusión. Para verificar cualquiera de estas
trayectorias hacen falta tres cosas que el apéndice \ref{cap:pedidos} detalla por organismo:

\begin{enumerate}[leftmargin=*]
\item \textbf{Partidas del Registro Civil}, que acreditan parentescos y descartan homonimias.
\item \textbf{Folios registrales}, que acreditan titularidades. El archivo entrega con precisión dos: el folio 342, asiento 438 del Libro B de La Caldera (finca Chalchanio), y el folio 374, asiento 475 del Libro D (finca Campo Alegre).
\item \textbf{Nóminas completas de funcionarios} del departamento por período, que para
1983--1987 hay que buscar como decretos de designación y no como proclamaciones de electos.
\end{enumerate}

Sin esas tres, ninguna de las coincidencias de esta lista es más que una coincidencia. Con
ellas, la tercera tesis de este libro se prueba o se cae.
